\scanpage{27}

\[
\left\{\begin{aligned}
g=0&\qquad\text{sss}\qquad n\leq5\\
g=1&\qquad\text{sss}\qquad n=6\\
g\geq2&\qquad\text{sss}\qquad n\geq7\;.
\end{aligned}\right.\tag*{(66)}
\]

Nous reviendrons de façon plus détaillée sur les cas où $g=0$ $(n\leq5)$ et $g=1$ $(n=6)$, par la suite. Pour l'instant, j'ai envie d'expliciter d'abord l\deleted{\unsure{e}}a fonctorialité des constructions faites, notamment pour $\Lambda$ variable, mais aussi pour $(M,e_M)$ variable.

$\overline{\underline{\mathrm{II}}}$

Revenons à la description combinatoire \unsure{standart|standard} des ``cartes triangulées'' orientées, via les ensembles à opérateurs sous le groupe cartographique triangulé \inserted{orienté (\unsure{pas précisé|pas pondéré|par prédisi|par prudence})}

\[
\mathbb{G}^{\unsure{o|0|\circ}}_{\underset{\deleted{\unsure{\infty}}}{3},\infty}=
\{\underset{\substack{\text{rotation}\\\text{autour}\\\text{d'un}\\\text{sommet}}}{\rho_0},
\underset{\substack{\text{rotation}\\\text{autour}\\\text{d'un}\\\text{centre}\\\text{d'arête}}}{\rho_1},
\underset{\substack{\text{rotation}\\\text{autour}\\\text{d'une face}}}{\rho_\infty}
\mid \rho_1^2=\rho_\infty^3=1,\deleted{\unsure{\}}}
\rho_\infty\rho_1\rho_0=1\}\tag*{(67)}
\]

D'autre part, on a un isomorphisme (plus ou moins)-canonique :

% archive sheet number: 26
